\documentclass[11pt]{article}
\usepackage[margin=1.1in]{geometry}
\usepackage{amsmath,amssymb,amsthm}
\usepackage{booktabs}
\usepackage[hidelinks]{hyperref}
\usepackage{microtype}

\newtheorem{corollary}{Corollary}
\setcounter{corollary}{-1} % so the first corollary is Corollary 0
\newcommand{\om}{\omega}
\newcommand{\vinf}{\|\om\|_\infty}
\newcommand{\Lam}{\Lambda}
\newcommand{\sL}{\sigma_{\!\Lambda}}

\title{Viscosity inverts the direction-regularity scaling of the
corner-flow blowup mechanism: a measured exponent crosses the type-I
exclusion threshold}
\author{Justin Hill\\ \small Independent Research \\ \small justin@epagoge.io}
\date{Draft 0.1 --- 2026-08-02\thanks{Every number in this paper
regenerates from bytes on disk with one command (Section~\ref{sec:repro}).
Claims are graded against a frozen standard; the grading ledger, the
confound engine, and the failed versions of our own claims are part of the
repository, not part of a memory hole.}}

\begin{document}
\maketitle

\begin{abstract}
For vorticity $\om$ with direction $\xi=\om/|\om|$ and peak set $P(t)$, the
quantity $\Lam=\sup_P |\nabla\xi|\,|\om|^{-1/2}$ is scale invariant under
the Navier--Stokes rescaling, and its exponent
$\sL = d\ln\Lam / d\ln\vinf$ turns the direction-coherence regularity
criteria of Constantin--Fefferman and Beir\~ao da Veiga--Berselli into a
single measurable number. We measure it. On the inviscid Luo--Hou corner
flow, where blowup is proven in the Boussinesq analogue by Chen--Hou, the
instrument returns $\sL=+1.00\pm0.03$, anchored by an exact amplitude
symmetry (predicted $\Lam$ ratio $0.707107$, measured $0.706891$) and a
$2\times2$ off-ray grid factorial (spread $0.029$). Under free-slip
viscosity, $\sigma(A)$ is not one number: every viscous run crosses over
from inviscid-like scaling at low amplitude to deep depletion at high
amplitude; window-averaged slopes are regime mixtures and we retract our
own earlier quotes of that kind. Measured in the asymptotic window
(top half of the cross-grid amplitude overlap, symmetric midpoint cut,
pre-registered), the deep-collapse exponent is
$\sL(10^{-4})=-1.12/-1.24$ (grids $128\times384$ / $256\times768$, spread
$0.121$) and $\sL(10^{-3})=-1.25/-1.29$ (spread $0.032$), against a
window-matched inviscid $+0.589/+0.585$ (spread $0.005$). Viscosity does
not merely damp the corner mechanism's alignment failure: it inverts it,
by roughly $1.8$ powers of $\vinf$, to a value far below both thresholds
of the exclusion chain ($-1/6$ for onset of geometric depletion, $-1/2$
for exclusion of type-I blowup by geometry plus energy, under one
structural hypothesis with a measured constant, Section~4). As far as a
live literature kill-search reaches, no direction-regularity exponent has
previously been measured as a function of viscosity on any blowup
candidate. On the dynamical side, a Lyapunov-certified march in the
rescaled frame (definiteness by Cholesky only; no eigenvalue enters any
decision) shows the collapse profile is an attractor: perturbed orbits
settle with $\cos_{\mathrm{step}}=0.986$ and contraction rate
$0.270\pm0.003$ per rescaled time unit, unanimous across five independent
seeds at the largest basin-validated amplitude.
\end{abstract}

\section{Introduction}\label{sec:intro}

Two problems sit on one function. If a viscous singularity inheriting the
corner mechanism exists, it must keep $\sL(\nu)$ above $-1/6$ as the
collapse deepens. If viscosity forces $\sL$ below $-1/2$ in the
high-vorticity regime, the geometric route closes the mechanism without
touching the supercritical energy problem: geometry, not energy, kills the
singularity. Nobody had $\sL(\nu)$ that our kill-search could find
(Section~\ref{sec:kill}). This paper reports it, for the
mechanism class where the inviscid answer is proven, at two viscosities,
on two grids, with the disagreement between grids quantified and under a
bar.

What is not new here, stated first. Direction-coherence regularity
criteria originate with Constantin and Fefferman \cite{CF93} and were
sharpened to the $1/2$-H\"older class by Beir\~ao da Veiga and Berselli
\cite{BdVB02,BdVB09}: the pairing of $|\nabla\xi|$ with the $1/2$ power of
amplitude lives in their sharp exponent before it lives in our observable.
Qualitative exclusion of type-I blowup under direction continuity is
likewise prior art: Giga--Miura \cite{GM11} for the infinite-energy class,
and Barker--Prange \cite{BP20} with scale-invariant alignment estimates in
the half-space with no-slip boundary, which is the Luo--Hou geometry
itself. The identity behind our instrument check
(Corollary~\ref{cor:zero}) is Constantin's classical decomposition of the
vorticity equation into magnitude and direction \cite{Con94}. A 2026
analytical line by Gruji\'c \cite{Gru26} pursues logarithmic depletion
through direction classes; it defines no observable and measures nothing.

What is new, correspondingly narrowed. First, the observable program:
treating the scale-forced pairing as a measurable scalar with an exponent,
and calibrating the instrument on a case where blowup is proven
\cite{CH22,CH25}, so that its verdict means something where blowup is
unknown. Second, the numbers: $\sL=+1.00\pm0.03$ on the inviscid corner
mechanism, the first direction-regularity exponent measured on any blowup
candidate as far as our search reaches. Third, the viscous inversion:
$\sL(\nu)$ measured at two viscosities with cross-grid certification,
crossing both thresholds deep in the collapse. Fourth, the quantified
exclusion chain (Section~\ref{sec:chain}): the known qualitative principle
\cite{GM11,BP20} sharpened to an explicit threshold on the measured
exponent by feeding the energy dissipation budget through Constantin's
kernel split, with the structural hypothesis carrying a measured constant.

The laboratory operates under a discipline that we consider part of the
result (Section~\ref{sec:repro}): a backwards proof of fifteen checks that
failed its own authors until the fit protocol was encoded, a confound
engine enumerating the ways each claim could be false, exact-symmetry
anchors with zero free parameters, and a standing watcher that reruns the
certification every twenty minutes. Several of our own earlier numbers
died by these instruments and are retained in the repository as corpses
with names.

\section{The observable}\label{sec:obs}

For a Navier--Stokes (or Euler, or Boussinesq) vorticity field, let
$\xi=\om/|\om|$ and $P(t)=\{x : |\om(x,t)|\ge\tfrac12\vinf\}$, and define
\begin{equation}
\Lam(t)=\sup_{P(t)}|\nabla\xi|\,|\om|^{-1/2},
\qquad
\sL=\frac{d\ln\Lam}{d\ln\vinf}.
\end{equation}
Under $u\to s\,u(sx,s^2t)$ the field $|\om|$ carries $s^2$ and
$|\nabla\xi|$ carries $s$, so $\Lam$ is exactly scale invariant; the $1/2$
power is forced by the group, and $|\xi|=1$ makes the direction field
globally controlled for free. The same pairing appears, in criterion form,
as the $\beta=1/2$ sharpness in \cite{BdVB02,BdVB09}; the observable is
that criterion turned into a number a simulation can return.

\begin{corollary}[of Constantin's $|\om|$-equation \cite{Con94}; stated as
a corollary, used as an instrument check]\label{cor:zero}
The magnitude of vorticity obeys
$\partial_t|\om|+u\!\cdot\!\nabla|\om|
=\alpha|\om|+\nu\left(\Delta|\om|-|\om|\,|\nabla\xi|^2\right)$,
where $\alpha=\xi\!\cdot\! S\xi$ is the stretching rate. At any growing
spatial maximum of $|\om|$ the Laplacian term is nonpositive, hence
\begin{equation}
\nu\,|\nabla\xi|^2\;\le\;\alpha \qquad\text{at the maximum.}
\end{equation}
\end{corollary}

This is a one-step corollary of a classical identity and we claim no
novelty for it; a literature kill-search retracted our earlier framing of
it as a theorem (Section~\ref{sec:kill}). It earns its place as a
validation: our data satisfies it with three orders of magnitude of margin
at every gated snapshot (campaign ledger), which any correct
direction-field computation must.

\section{Instruments}\label{sec:instruments}

The DNS side is a Dedalus-based axisymmetric Boussinesq solver in the
Luo--Hou corner scenario, with free-slip viscosity in 5D-Laplacian form
($\nu\Delta_5$ on $u_1$ and $\om_1$, $\Delta_5=\partial_{rr}
+(3/r)\partial_r+\partial_{zz}$), implicit in the IMEX stepper, and a
quartic initial condition exactly compatible with the Neumann wall
condition. Trust gates per snapshot: spectral tail at or below $10^{-6}$
and signed drift of the conserved circulation invariant $\sup|r^2u_1|$
(decay under viscosity is physics; growth is a violation). $\Lam$ is evaluated on the $2\times$-HWHM box at the vorticity
peak; runs are trimmed at the first $10\times$ amplitude jump. The
exponent is the slope of $\ln\Lam$ against $\ln\vinf$ over gated
snapshots, with pair-resampled bootstrap confidence intervals ($10^4$
resamples) everywhere a slope is quoted.

The rescaled-frame side is a corner-frame $s$-march on the certified
self-similar profile: backward Euler on the index-2 DAE with the algebraic
manifold re-imposed at every iterate, a frozen reduced Jacobian whose
exact validity on $\ker(C_g)$ is a closed-form identity (verified at
$4.6\times10^{-14}$ against the exact marcher before being trusted), and a
Lyapunov $P$-norm gate in which definiteness is decided by Cholesky
factorization only. No eigenvalue of any operator enters any decision
anywhere in this paper; that refusal is standing and was never violated
under pressure.

\section{The exclusion chain, quantified}\label{sec:chain}

Constantin's identity for the stretching rate $\alpha=\xi\!\cdot\! S\xi$
consumes the Biot--Savart kernel geometry: misalignment is the only way to
stretch. Splitting the integral at radius $\rho$, the near field is
controlled by $\Lam\,\vinf^{3/2}\rho$ and the far field by the energy
through Cauchy--Schwarz; optimizing $\rho$,
\begin{equation}
\alpha \;\lesssim\;
\left(\Lam\,\vinf^{3/2}\right)^{3/5}\|\om\|_{L^2}^{2/5}.
\end{equation}
Hypothesis (H): the peak set is a single coherent structure on which the
gradient bound holds out to the optimal $\rho$. The hypothesis carries a
measured constant ($\lambda_0=5.0$) and the chain inequality held at 21
of 21 gated snapshot pairs with the stretching rate inferred independently
(campaign ledger, STANDARD.md; that run predates the artifacts retained in
this campaign's outputs). Feeding the only global control Navier--Stokes provides
(the energy dissipation identity) through this bound, with blowup rate
$\vinf\sim(T-t)^{-\gamma}$, $\gamma\ge1$:
\begin{align}
\sL\le-\tfrac16 &: \text{geometric depletion begins to act at all,}\\
\sL<-\tfrac12 &: \text{type-I blowup excluded by geometry and energy alone.}
\end{align}
The qualitative principle here is not ours: under a type-I condition,
direction continuity where vorticity is large already excludes blowup
\cite{GM11}, including in the half-space with no-slip boundary
\cite{BP20}. What the chain above adds is the exponent: an explicit
threshold on a measurable number, so that a simulation, or eventually an
experiment, can locate a given flow on the right or wrong side of it.

\section{The inviscid calibration}\label{sec:inviscid}

On the inviscid corner flow the instrument returns
\begin{equation}
\sL = +1.00 \pm 0.03.
\end{equation}
Two validations with zero free parameters between them. First, the exact
amplitude symmetry $u\to\lambda u$, $t\to t/\lambda$ of the Euler
equations predicts a $\Lam$ ratio of $1/\sqrt2=0.707107$ between paired
runs; measured $0.706891\pm0.001441$, three parts in ten thousand, with
the predicted $|\nabla\xi|$ ratio of $1$ measured at $0.999746$. Second, a
$2\times2$ off-ray factorial refining $N_z$ and $N_r$ independently
returns slopes $+1.0159$, $+0.9873$, $+1.0027$, $+0.9978$ (spread $0.029$,
main effects at or below $0.017$), rejecting the axial-mesh and
single-ray confounds by test rather than by assertion.

Reading: $|\nabla\xi|\sim\vinf^{3/2}$. Direction regularity fails a full
$3/2$ of a power above the exclusion threshold; the Euler mechanism blows
up because its geometry is maximally non-depleting. The observable returns
the correct sign on a case where blowup is proven in the Boussinesq
analogue \cite{CH22,CH25}, which is precisely what licenses pointing it at
cases where the answer is unknown.

Refinement recorded during the viscous campaign: the full-window $+1.00$
is amplitude-composite. Fit on matched cross-grid amplitude windows, the
inviscid slope runs $+1.42/+1.36$ in the lower half and $+0.589/+0.585$
in the deep half (cross-grid spread $0.005$, the tightest certification in
this paper; retained artifact \texttt{M4\_SIGMA\_DEEP\_NU0.out}). The symmetry and factorial validations certify the instrument, not
slope constancy across amplitude; window-matched comparisons below
therefore use $+0.59$ as the inviscid contrast value.

\section{The viscous measurement}\label{sec:viscous}

Runs at $\nu\in\{10^{-4},10^{-3}\}$ on grids $128\times384$ and
$256\times768$ with inviscid references on both grids. Three findings, in
discovery order.

(i) \emph{The common-amplitude overlap window reproduces the calibration.}
Fit on the same amplitude range, the two inviscid grids give
$+1.016/+1.033$ (spread $0.017$) against the calibrated $+1.00\pm0.03$.
The cross-grid disagreement of naive full-window fits ($0.147$) was a
window artifact, not a physics disagreement.

(ii) \emph{At finite $\nu$, $\sigma(A)$ is not one number.} Every viscous
run crosses over from inviscid-like slope ($+0.39$ to $+0.81$) at low
amplitude to deep depletion near $-1.2$ at high amplitude, on both grids,
at both viscosities. A window-averaged slope is a regime mixture; our own
earlier quotes of that kind (window averages near $-0.5$, and a
single-grid ladder $-0.37$ to $-0.57$) are retracted as measurements of a
mixture, and the retraction is part of the repository record.

(iii) \emph{Measured where the asymptotic regime actually lives} (top half
of the cross-grid overlap in $\ln A$, symmetric midpoint cut,
pre-registered before the cut-sensitivity scan), the deep-collapse
exponent is
\begin{center}
\begin{tabular}{lccc}
\toprule
 & $128\times384$ & $256\times768$ & spread \\
\midrule
$\nu=10^{-4}$ & $-1.116$ & $-1.237$ & $0.121$ \\
$\nu=10^{-3}$ & $-1.254$ & $-1.285$ & $0.032$ \\
\midrule
$\nu=0$ (matched window) & $+0.589$ & $+0.585$ & $0.005$ \\
\bottomrule
\end{tabular}
\end{center}
with bootstrap 95\% intervals of half-width $\pm0.14$ to $\pm0.30$ on each fit
and cut-position sensitivity reported at $0.4/0.5/0.6$ of the overlap
span: the magnitude sits between $-0.99$ and $-1.31$ at every cut, on every
grid, at both viscosities. Against the window-matched inviscid $+0.59$,
viscosity inverts the direction-regularity scaling by roughly $1.8$ powers
of $\vinf$, to a value below both thresholds of Section~\ref{sec:chain}.

Honest bounds on what this shows. Two viscosities, one mechanism class,
one scenario geometry, laptop resolution; the deep windows carry 11 to 14
snapshots each; the crossover amplitude $A_c(\nu)$ is visible but not yet
resolved as its own observable; and the exclusion chain runs through
Hypothesis (H) with its measured constant. What the measurement supports,
within those bounds, is exact: on the strongest known blowup mechanism of
this class, the geometric quantity that must stay above $-1/6$ for the
mechanism to survive viscosity is measured at $-1.2$ and below, deep in
the collapse, grid-certified, at both viscosities tried.

\section{The kill-search}\label{sec:kill}

A live literature search (2026-08-02; arXiv, Springer, journal indices;
queries on direction-coherence criteria, H\"older-$1/2$ direction results,
scale-invariant alignment estimates, type-I exclusion via direction,
Luo--Hou alignment measurements, and name collisions for $\sL$) graded
every novelty-bearing claim. One claim died: our Theorem~0 is a one-step
corollary of Constantin's $|\om|$-equation and is demoted to
Corollary~\ref{cor:zero} above. Theorem-1-as-principle died too, as it
should: the qualitative type-I exclusion belongs to \cite{GM11} and
\cite{BP20}; what survives is the quantification. The observable program,
the inviscid measurement, and the viscous inversion survived the search
outright, with the standing caveat that survived means \emph{not found by
this search}, not \emph{proven absent}. Closest adjacent art in the
numerical literature: Hou's dynamic-depletion diagnostics (alignment with
strain eigenvectors, vortex-line regularity; qualitative)
\cite{Hou09,LH14}, and Kerr's numerical tests of direction-derived
regularity conditions on antiparallel vortex tubes \cite{Kerr13}.

\section{Reproducibility and discipline}\label{sec:repro}

\begin{verbatim}
python verify.py        # the backwards proof: 15 checks, 3 tiers
python mythos.py        # the confound engine
./done.sh               # the milestone ladder, M1-M5, all PASS
cat VERIFY_STATUS       # what the standing watcher currently holds
\end{verbatim}

\texttt{verify.py} holds at 15/15 with EXACT-tier worst deviation
$3.05\times10^{-4}$ against references forced by the equations before any
measurement existed; input checksums are in the verify output; a watcher
reruns the full certification every twenty minutes and files a regression
report if anything drifts. The laboratory was built by a human directing
an AI system; the AI invented a branch that did not exist, fit through
windows that moved, and glossed a mechanism its own certification refuted
within the hour. Every failure was caught the same day by exact symmetry
anchors, context-free review, or the backwards proof, and the corpses are
kept on the page. The discipline that survives contact with an AI that
launders inference is part of what this repository demonstrates.

\section{Open}\label{sec:open}

The strip $-1/2\le\sigma\le0$: Corollary~\ref{cor:zero} forbids the top
under viscosity at a growing maximum; the chain of
Section~\ref{sec:chain} needs the bottom. Close the strip analytically and
type-I Navier--Stokes blowup is finished for this route. The crossover
amplitude $A_c(\nu)$ and the $\sL(\nu)$ curve beyond the two certified
viscosities are unmeasured (stream data exists at
$\nu=3\times10^{-4},3\times10^{-5},10^{-5}$, extending coverage one
further decade; snapshots were not retained).
The wandering sector remains the measured frontier: generic near-wall data
dies too young at laptop resolution to read, and that failure is the
recorded reason the sector is unexplored.

\begin{thebibliography}{99}

\bibitem{CF93} P.~Constantin, C.~Fefferman.
Direction of vorticity and the problem of global regularity for the
Navier--Stokes equations.
\emph{Indiana Univ. Math. J.} \textbf{42}, 775--789 (1993).

\bibitem{Con94} P.~Constantin.
Geometric statistics in turbulence.
\emph{SIAM Review} \textbf{36}, 73--98 (1994).

\bibitem{BdVB02} H.~Beir\~ao da Veiga, L.C.~Berselli.
On the regularizing effect of the vorticity direction in incompressible
viscous flows.
\emph{Differential Integral Equations} \textbf{15} (2002).

\bibitem{BdVB09} H.~Beir\~ao da Veiga, L.C.~Berselli.
Navier--Stokes equations: Green's matrices, vorticity direction, and
regularity up to the boundary (2009).

\bibitem{GM11} Y.~Giga, H.~Miura.
On vorticity directions near singularities for the Navier--Stokes flows
with infinite energy.
\emph{Comm. Math. Phys.} (2011). doi:10.1007/s00220-011-1197-x

\bibitem{BP20} T.~Barker, C.~Prange.
Scale-invariant estimates and vorticity alignment for Navier--Stokes in
the half-space with no-slip boundary conditions.
\emph{Arch. Ration. Mech. Anal.} \textbf{235}, 881--926 (2020).
arXiv:1906.08225.

\bibitem{Gru26} Z.~Gruji\'c.
Logarithmic depletion of vortex stretching and singularity evasion in the
3D Navier--Stokes equations.
arXiv:2607.08866 (2026).

\bibitem{Hou09} T.Y.~Hou.
Blow-up or no blow-up? A unified computational and analytic approach to 3D
incompressible Euler and Navier--Stokes equations.
\emph{Acta Numerica} (2009).

\bibitem{LH14} G.~Luo, T.Y.~Hou.
Toward the finite-time blowup of the 3D axisymmetric Euler equations: a
numerical investigation.
\emph{Multiscale Model. Simul.} (2014).

\bibitem{CH22} J.~Chen, T.Y.~Hou.
Stable nearly self-similar blowup of the 2D Boussinesq and 3D Euler
equations with smooth data I: Analysis.
arXiv:2210.07191.

\bibitem{CH25} J.~Chen, T.Y.~Hou.
Stable nearly self-similar blowup of the 2D Boussinesq and 3D Euler
equations with smooth data II: Rigorous Numerics.
\emph{Multiscale Model. Simul.} \textbf{23}(1), 25--130 (2025).
arXiv:2305.05660.

\bibitem{Kerr13} R.M.~Kerr.
Bounds for Euler from vorticity moments and line divergence.
arXiv:1212.1106.

\bibitem{Surv21} (Survey.)
A survey of geometric constraints on the blowup of solutions of the
Navier--Stokes equation.
arXiv:2111.00040.

\end{thebibliography}

\end{document}
